\section*{Part II: Cost Behavior and the Economics of Production}

\subsection*{8. Fixed, Variable, Step-Fixed, and Mixed Costs}

A \textbf{fixed cost} remains approximately constant over a defined range of output and time: a base software license, minimum rent, or a salaried program manager. A \textbf{variable cost} increases with units: payment processing, material inputs, per-task inference, or reviewer time that scales with submissions. A \textbf{step-fixed cost} stays level until demand requires another server tier, teaching assistant, machine, classroom, or shift. A \textbf{mixed cost} has a base fee plus usage or overtime components.

Costs must be classified relative to a specific decision horizon. Salaries may be fixed for a month but adjustable over a year. Cloud infrastructure may be variable over daily workloads yet contractually committed for an annual period. A rented room may be avoidable for a pilot but fixed after signing a lease. Therefore ``fixed'' is not an eternal property of an expense; it is a description of how that expense behaves within an operating decision.

The founder should calculate contribution at realistic capacity bands, not assume a smooth cost curve. An attractive cohort model at 18 learners may become unattractive at 21 when another qualified reviewer is needed.

\subsection*{9. Direct Cost Is Not Identical to Variable Cost}

A direct cost can be traced to a product or contract: a named instructor, project travel, or specific cloud workload. It might still be fixed within the contract. Variable cost changes with units but may be difficult to trace to a single customer. These accounting and managerial categories answer different questions.

For a learning cohort, the lead instructor's 30 hours are directly attributable to the cohort but may not change whether 12 or 15 students enroll. A reviewer paid per submission scales with completed work. Marketing allocated to the cohort could be direct while its cost per enrollment depends on conversion. Distinguish attribution from cost behavior or contribution calculations will become inconsistent.

\subsection*{10. Marginal Cost and the Point Where Scale Stops Helping}

Marginal cost is the additional cost of providing 1 more unit at a given operating state. A copied recording may cost nearly nothing to distribute, but an additional enrolled learner can increase assessment, feedback, support, refunds, compute, and administrative effort. The relevant marginal cost is the \textbf{whole additional promise}, not a cheap technical component.

There may be unused capacity that makes the next customer inexpensive, followed by a threshold where a new staff member or facility is required. A founder should model both the incremental cost \emph{within} the current capacity band and the new cost \emph{after} the threshold. Otherwise optimistic average costs conceal a future staffing cliff.

\subsection*{11. Average Cost and Economies of Scale}

Average cost divides total cost by output:
\begin{equation}
AC(Q) = \frac{TC(Q)}{Q}.
\end{equation}
When fixed costs are spread across more units, average cost can decline. Standardization, procurement scale, reusable intellectual property, learning curves, and better utilization can reinforce the effect. But scale also brings coordination, quality control, management, churn, compliance, and service complexity.

A founder should not assume that every service inherits software-like economics. In high-touch education, more learners may require nearly proportional skilled review, and raising utilization too far can worsen latency and outcomes. True scale advantage comes when \textbf{total cost grows more slowly than valuable delivered output}, without hiding quality erosion.

\subsection*{12. Economies of Scope and Reusable Assets}

A firm enjoys economies of scope when a shared capability lowers the combined cost of producing related offerings relative to producing them separately. A rigorous assessment framework, reusable code examples, learning platform, or trusted technical faculty might support D2C, college, and enterprise formats. But the services may also require different buyer systems, compliance, schedules, and customer support.

Reusability should be measured: What fraction of content, technology, assessment, distribution, or delivery actually transfers? How much adaptation is needed? What conflicts arise? If serving an enterprise forces the teaching team to abandon scheduled cohort feedback, the apparent scope advantage may be overwhelmed by opportunity cost and quality damage.
